\item Una cámara de combustión está a $750\text{ K}$ y la temperatura ambiente a $300\text{ K}$. La eficiencia de una máquina de Carnot que hace $150\text{ J}$ de trabajo mientras transporta energía entre estos baños de temperaturas constantes es 60.0\%. La máquina de Carnot debe tomar energía a $150\text{ J}/0.600 = 250\text{ J}$ del depósito caliente y colocar $100\text{ J}$ de energía por calor en el ambiente. Para seguir el razonamiento de Carnot, suponga que alguna otra máquina térmica S podría tener una eficiencia de 70.0\%. 

(a) Encuentre la entrada de energía y la salida de energía desperdiciada de la máquina S conforme consume $150\text{ J}$ de trabajo. 
(b) Deje que la máquina S funcione como en el inciso (a) y opere la máquina de Carnot de reversa entre los mismos depósitos. El trabajo de salida de la máquina S es el trabajo de entrada para el refrigerador de Carnot. Encuentre la energía total transferida a o desde la cámara de combustión y la energía total transferida a o desde el ambiente cuando ambas máquinas funcionan juntas. 
(c) Explique cómo los resultados de los incisos (a) y (b) demuestran que se viola el enunciado de Clausius de la segunda ley de la termodinámica. 
(d) Encuentre la entrada de energía y la salida de trabajo de la máquina S conforme saca energía de escape de $100\text{ J}$. Deje que la máquina S opere como en el inciso (c) y aporte $150\text{ J}$ de su trabajo de salida para que funcione la máquina de Carnot en reversa. Encuentre 
(e) la energía total que la cámara proporciona cuando ambas máquinas funcionan juntas, 
(f) el trabajo total de salida y 
(g) la energía total transferida al ambiente. 
(h) Explique cómo los resultados demuestran que se viola el enunciado de Kelvin-Planck de la segunda ley. Por tanto, la suposición acerca de la eficiencia de la máquina S debe ser falsa. 
(i) Deje que las máquinas operen juntas a través de un ciclo, como en el inciso (d). Encuentre el cambio en entropía del Universo. 
(j) Explique cómo el resultado prueba que se viola el enunciado de entropía de la segunda ley.
